\section{Statements and Operators}

\subsection{Section Overview}
In this section, we will delve into expressions, statements, and the various operators C++ offers. Operators are the symbols that perform operations on data. We will cover assignment, arithmetic, logical, relational, and other important operators. Understanding these is key to performing calculations and making decisions in your programs.

\subsection{Expressions and Statements}
An \textbf{expression} is a sequence of operators and operands that specifies a computation. An expression has a value. For example, \texttt{2 + 3} is an expression whose value is 5.

A \textbf{statement} is a complete instruction to the computer. In C++, most statements are terminated with a semicolon. A statement may or may not have a value. For example, \texttt{int x = 5;} is a statement.

\subsection{Using Operators}
Operators are special symbols that perform operations on one or more operands.
\begin{itemize}
	\item \textbf{Unary} operators act on one operand (e.g., \texttt{-x}).
	\item \textbf{Binary} operators act on two operands (e.g., \texttt{x + y}).
	\item \textbf{Ternary} operators act on three operands (e.g., \texttt{x > y ? 1 : 0}).
\end{itemize}

\subsection{The Assignment Operator}
The assignment operator (\texttt{=}) assigns the value on its right to the variable on its left.
\begin{lstlisting}
	int x;
	x = 5; // x is now 5

	int y = 10; // initialization is not assignment
	y = x;      // y is now 5
\end{lstlisting}
The assignment operator is a binary operator.

\subsection{Arithmetic Operators}
C++ provides the following arithmetic operators:
\begin{itemize}
	\item \texttt{+} addition
	\item \texttt{-} subtraction
	\item \texttt{*} multiplication
	\item \texttt{/} division
	\item \texttt{\%} modulus (remainder)
\end{itemize}
\begin{lstlisting}
	int a = 10;
	int b = 3;

	int sum = a + b; // 13
	int diff = a - b; // 7
	int prod = a * b; // 30
	int quot = a / b; // 3 (integer division)
	int rem = a % b;  // 1

	double c = 10.0;
	double d = 3.0;
	double quot_double = c / d; // 3.333...
\end{lstlisting}

\subsection{Coding Exercise 7: Using the Assignment Operator}
Declare two integer variables, \texttt{a} and \texttt{b}. Initialize \texttt{a} to 10. Use the assignment operator to make \texttt{b} equal to \texttt{a}. Print the value of \texttt{b}.
\begin{lstlisting}
	#include <iostream>

	int main() {
		int a = 10;
		int b;
		b = a;
		std::cout << "The value of b is: " << b << std::endl;
		return 0;
	}
\end{lstlisting}

\subsection{Coding Exercise 8: Using the Arithmetic Operators}
Write a program that asks the user to enter two integers, then prints their sum, difference, product, and quotient (integer division).
\begin{lstlisting}
	#include <iostream>

	int main() {
		int num1, num2;
		std::cout << "Enter two integers: ";
		std::cin >> num1 >> num2;
		std::cout << "Sum: " << num1 + num2 << std::endl;
		std::cout << "Difference: " << num1 - num2 << std::endl;
		std::cout << "Product: " << num1 * num2 << std::endl;
		std::cout << "Quotient: " << num1 / num2 << std::endl;
		return 0;
	}
\end{lstlisting}

\begin{remark}
	This program has a hidden hazard: if the user enters \texttt{0} for \texttt{num2}, the division \texttt{num1 / num2} is \textbf{integer division by zero}—this is undefined behavior and will typically crash the program. In production code, always guard against zero before dividing:
	\begin{lstlisting}
		if (num2 != 0)
			std::cout << "Quotient: " << num1 / num2 << std::endl;
		else
			std::cout << "Cannot divide by zero." << std::endl;
	\end{lstlisting}
\end{remark}

\subsection{Increment and Decrement Operators}
The increment (\texttt{++}) and decrement (\texttt{-{}-}) operators increase or decrease a variable's value by one. They can be used in prefix or postfix form.
\begin{lstlisting}
	int counter = 10;

	// Postfix increment
	std::cout << counter++ << std::endl; // prints 10, counter becomes 11

	// Prefix increment
	std::cout << ++counter << std::endl; // counter becomes 12, prints 12

	// Postfix decrement
	std::cout << counter-- << std::endl; // prints 12, counter becomes 11

	// Prefix decrement
	std::cout << --counter << std::endl; // counter becomes 10, prints 10
\end{lstlisting}

\subsection{Mixed Expressions and Conversions}
When you mix types in an expression, C++ performs conversions. Generally, the ``lower'' type is promoted to the ``higher'' type. For example, in an operation involving an \texttt{int} and a \texttt{double}, the \texttt{int} is converted to a \texttt{double}.
\begin{lstlisting}
	int total {100};
	int count {8};
	double average;

	average = total / count; // integer division, average is 12.0

	average = static_cast<double>(total) / count; // correct: 12.5
\end{lstlisting}

\begin{remark}
	The first line is a subtle trap: \texttt{total / count} is evaluated as integer division (\texttt{100 / 8 = 12}) \textbf{before} the result is stored in \texttt{average}. The truncation to 12 happens at the division step, not when assigning to \texttt{double}. Casting \textit{either} operand to \texttt{double} forces the entire division to use floating-point arithmetic.
\end{remark}

\subsection{Testing for Equality}
The equality operator (\texttt{==}) tests if two values are equal. The inequality operator (\texttt{!=}) tests if two values are not equal.
\begin{lstlisting}
	int a = 10, b = 10, c = 20;

	bool equal = (a == b); // true
	bool not_equal = (a != c); // true
\end{lstlisting}
\begin{remark}
	Do not confuse the equality operator (\texttt{==}) with the assignment operator (\texttt{=}). This is a very common mistake.
\end{remark}

\subsection{Relational Operators}
Relational operators are used to compare values.
\begin{itemize}
	\item \texttt{>} greater than
	\item \texttt{<} less than
	\item \texttt{>=} greater than or equal to
	\item \texttt{<=} less than or equal to
\end{itemize}
These operators return \texttt{true} or \texttt{false}.
\begin{lstlisting}
	int a = 10, b = 20;
	bool is_greater = (b > a); // true
	bool is_less_equal = (a <= b); // true
\end{lstlisting}

\subsection{Logical Operators}
Logical operators are used to combine or modify logical values (\texttt{true} or \texttt{false}).
\begin{itemize}
	\item \texttt{\&\&} AND: \texttt{true} if both operands are \texttt{true}.
	\item \texttt{||} OR: \texttt{true} if at least one operand is \texttt{true}.
	\item \texttt{!} NOT: inverts the truth value.
\end{itemize}
\begin{lstlisting}
	bool a = true, b = false;
	bool both_true = a && b; // false
	bool at_least_one_true = a || b; // true
	bool not_a = !a; // false
\end{lstlisting}

\begin{remark}
	\textbf{Short-circuit evaluation}: \texttt{\&\&} and \texttt{||} stop evaluating as soon as the result is determined. For \texttt{A \&\& B}, if \texttt{A} is \texttt{false}, \texttt{B} is never evaluated. For \texttt{A || B}, if \texttt{A} is \texttt{true}, \texttt{B} is never evaluated. This matters when the right operand has a side effect or can crash:
	\begin{lstlisting}
		// Safe: denominator is checked before division
		if (count != 0 && total / count > 10) { ... }

		// Safe: ptr is checked before dereferencing
		if (ptr != nullptr && ptr->value > 0) { ... }
	\end{lstlisting}
\end{remark}

\subsection{Compound Assignment Operators}
Compound assignment operators combine an arithmetic operation with assignment.
\begin{itemize}
	\item \texttt{+=} e.g., \texttt{a += 5} is \texttt{a = a + 5}
	\item \texttt{-=} e.g., \texttt{a -= 5} is \texttt{a = a - 5}
	\item \texttt{*=} e.g., \texttt{a *= 5} is \texttt{a = a * 5}
	\item \texttt{/=} e.g., \texttt{a /= 5} is \texttt{a = a / 5}
	\item \texttt{\%=} e.g., \texttt{a \%= 5} is \texttt{a = a \% 5}
\end{itemize}
\begin{lstlisting}
	int x = 10;
	x += 5; // x is now 15
	x *= 2; // x is now 30
\end{lstlisting}

\subsection{Operator Precedence}
Operator precedence determines the order in which operators in an expression are evaluated. For example, \texttt{*} and \texttt{/} have higher precedence than \texttt{+} and \texttt{-}. You can use parentheses \texttt{()} to override the default precedence. A full precedence table can be found at \href{https://en.cppreference.com/w/cpp/language/operator_precedence}{cppreference.com}. When in doubt, use parentheses.

\begin{example}[Precedence Pitfalls]{\leavevmode}
	\begin{lstlisting}
		// Without parentheses - precedence may surprise you
		bool result = 5 + 3 > 6 && 2 * 4 == 8;
		// Evaluated as: ((5+3) > 6) && ((2*4) == 8) => true && true => true

		// Confusing: assignment has very low precedence
		bool flag;
		flag = 3 > 2;   // flag = (3 > 2) = true  -- correct
		// flag = 3 > 2 == 1;  // flag = (3 > (2 == 1)) -- NOT what you want

		// Rule of thumb: use parentheses whenever mixing
		// arithmetic, relational, and logical operators
		bool can_proceed = (age >= 18) && (score > 50);
	\end{lstlisting}
\end{example}

\subsection{Coding Exercise 9: Logical Operators and Operator Precedence --- Can you work?}
Write a program that determines if a person can work. A person can work if they are between 18 and 65 years old (inclusive). Ask the user for their age.
\begin{lstlisting}
	#include <iostream>

	int main() {
		int age;
		std::cout << "Enter your age: ";
		std::cin >> age;
		bool can_work = (age >= 18) && (age <= 65);
		// std::boolalpha makes cout print "true"/"false"
		// instead of 1/0
		std::cout << "Can you work? "
			<< std::boolalpha << can_work << std::endl;
		return 0;
	}
\end{lstlisting}

\subsection{Section Challenge}
Write a program that asks the user to enter a three-digit integer. The program should then print the digits of the integer separately. For example, if the user enters 123, the program should print 1, then 2, then 3.

\textit{Hint: use the division and modulus operators.}

\subsection{Section Challenge --- Solution}
\begin{lstlisting}
	#include <iostream>

	int main() {
		int number;
		std::cout << "Enter a three-digit number: ";
		std::cin >> number;

		int digit1 = number / 100;
		int digit2 = (number / 10) % 10;
		int digit3 = number % 10;

		std::cout << "The digits are: "
			<< digit1 << ", " << digit2 << ", " << digit3 << std::endl;

		return 0;
	}
\end{lstlisting}

\subsection{Quiz 5: Section 08 Quiz}
\begin{enumerate}
	\item Which of the following is the equality operator?
	\begin{enumerate}
		\item[a)] \texttt{=}
		\item[b)] \texttt{:=}
		\item[c)] \texttt{==}
		\item[d)] \texttt{eq}
	\end{enumerate}

	\item What is the value of the expression \texttt{10 \% 3}?
	\begin{enumerate}
		\item[a)] 3
		\item[b)] 1
		\item[c)] 0
		\item[d)] 3.33
	\end{enumerate}

	\item What does the \texttt{\&\&} operator do?
	\begin{enumerate}
		\item[a)] Logical OR
		\item[b)] Logical NOT
		\item[c)] Bitwise AND
		\item[d)] Logical AND
	\end{enumerate}
\end{enumerate}

\textbf{Answers:} 1.\ (c), 2.\ (b), 3.\ (d)
